\documentclass[aspectratio=169, 10pt]{beamer}
\usetheme{metropolis}

% --- Edinburgh colours ----------------------------------------
\definecolor{UoERed}{HTML}{7A2318}
\definecolor{UoEGold}{HTML}{B8860B}
\definecolor{CodeBG}{HTML}{F7F3ED}
\definecolor{CodeFrame}{HTML}{D5C9B1}

\setbeamercolor{palette primary}{bg=UoERed, fg=white}
\setbeamercolor{frametitle}{bg=UoERed, fg=white}
\setbeamercolor{title separator}{fg=UoEGold}
\setbeamercolor{progress bar}{fg=UoEGold, bg=UoEGold!20}
\setbeamercolor{alerted text}{fg=UoERed}
\setbeamercolor{block title}{bg=UoERed!15, fg=UoERed}
\setbeamercolor{block body}{bg=UoERed!5}

% --- Packages -------------------------------------------------
\usepackage[T1]{fontenc}
\usepackage{lmodern}
\usepackage{amsmath, amssymb, mathtools}
\usepackage{listings}
\usepackage{graphicx, booktabs}
\usepackage{tikz}
\usetikzlibrary{arrows.meta, positioning, calc, decorations.pathreplacing}
\usepackage{hyperref}
\hypersetup{colorlinks=true, linkcolor=UoERed, urlcolor=UoERed!70!black}
\setbeamertemplate{navigation symbols}{}

\lstset{
  language=Python,
  basicstyle=\ttfamily\scriptsize,
  keywordstyle=\color{UoERed}\bfseries,
  commentstyle=\color{gray}\itshape,
  stringstyle=\color{UoEGold},
  backgroundcolor=\color{CodeBG},
  frame=single,
  rulecolor=\color{CodeFrame},
  numbers=none, breaklines=true, showstringspaces=false, tabsize=4,
  xleftmargin=4pt, xrightmargin=4pt, aboveskip=6pt, belowskip=4pt,
}

\AtBeginSection[]{
  \begin{frame}
    \vfill\centering
    \usebeamerfont{title}\insertsectionhead\par
    \vfill
  \end{frame}
}

\title{Week 4 --- Deep Equilibrium Networks I}
\subtitle{The DEQN Principle \& the Brock--Mirman Growth Model}
\author{Dr Juan Zurita}
\institute{Deep Learning for Macroeconomics\\
The University of Edinburgh~$\cdot$~School of Economics}
\date{}

\begin{document}
\maketitle

\begin{frame}{From VFI to DEQNs}
In PNM you solved the growth model by \textbf{value-function iteration} on a grid.

\bigskip
\textbf{DEQN idea:} parameterise the policy function as a neural network and train it to satisfy the \alert{equilibrium conditions}.

\bigskip
\begin{block}{Azinovic, Gaegauf \& Scheidegger (2022)}
``Deep Equilibrium Nets'' --- \textit{International Economic Review}
\end{block}
\end{frame}

\begin{frame}{The DEQN Recipe}
\begin{enumerate}
\item \textbf{Define} the economic model (states, controls, equilibrium conditions)
\item \textbf{Parameterise} policy functions: $c = \mathcal{N}_{\boldsymbol{\theta}}(k)$
\item \textbf{Construct loss}: squared Euler-equation residuals
\item \textbf{Train}: minimise the loss with Adam
\item \textbf{Verify}: compare with known solutions or refined grids
\end{enumerate}

\bigskip
No grid, no discretisation of the state space --- just sample points.
\end{frame}

\begin{frame}{Brock--Mirman Model}
$\max \sum_{t=0}^{\infty} \beta^t \ln c_t$ \quad s.t.\ $k_{t+1} = k_t^\alpha - c_t$

\bigskip
Euler equation: $\frac{1}{c_t} = \beta\,\frac{\alpha k_{t+1}^{\alpha-1}}{c_{t+1}}$

\bigskip
\begin{block}{Why this model?}
It has a \textbf{closed-form solution}: $c_t = (1 - \alpha\beta)\,k_t^\alpha$.

We can verify the neural network against the true answer.
\end{block}
\end{frame}

\begin{frame}{Loss Function Design}
The Euler-equation residual at a sample point $k$:

$$R(k; \boldsymbol{\theta}) = 1 - \beta\,\frac{\alpha k'^{\alpha-1} \cdot c(k; \boldsymbol{\theta})}{c(k'; \boldsymbol{\theta})}$$

where $k' = k^\alpha - c(k; \boldsymbol{\theta})$.

\bigskip
Loss: $\mathcal{L}(\boldsymbol{\theta}) = \frac{1}{N}\sum_{i=1}^{N} R(k_i; \boldsymbol{\theta})^2$

\bigskip
Train until $\mathcal{L} < 10^{-8}$ --- machine-precision accuracy.
\end{frame}

\begin{frame}{Summary}
\begin{itemize}
\item DEQNs replace grid-based VFI with neural-network policy functions
\item The loss function encodes economic equilibrium conditions
\item Brock--Mirman provides a verifiable first test case
\item The method is grid-free and scales to higher dimensions
\end{itemize}

\bigskip
\textbf{Next week:} adding uncertainty --- stochastic models with Gauss--Hermite quadrature.
\end{frame}

\end{document}
